\documentclass[tikz,border=5mm]{standalone}
\usepackage{luatexja}
\usepackage{tikz}
\usetikzlibrary{arrows.meta, positioning, fit, backgrounds, calc, shapes.geometric}

\begin{document}

\tikzset{
    block/.style={
        rectangle,
        draw=black,
        thick,
        fill=#1,
        minimum width=2.8cm,
        minimum height=1cm,
        font=\small,
        text centered,
        rounded corners=3pt
    },
    block/.default=blue!15,
    datablock/.style={
        block=green!15,
        minimum width=2.2cm
    },
    modelblock/.style={
        block=orange!20,
        minimum width=2.5cm
    },
    lossblock/.style={
        block=red!15,
        minimum width=2.2cm
    },
    arrow/.style={
        -Stealth,
        thick
    },
    dashedarrow/.style={
        -Stealth,
        thick,
        dashed
    },
    loopbox/.style={
        rectangle,
        draw=black!60,
        thick,
        rounded corners=5pt,
        fill=gray!5
    },
    looplabel/.style={
        rectangle,
        fill=gray!30,
        font=\footnotesize\bfseries,
        inner sep=3pt,
        rounded corners=2pt
    },
    note/.style={
        font=\footnotesize,
        text=gray!70
    }
}

\begin{tikzpicture}[node distance=1.2cm and 1.5cm]
% === データ ===
\node[block=green!15, minimum width=8cm] (databox) {データ};

% === 前処理 ===
\node[block=cyan!15, below=1.5cm of databox, minimum width=8cm] (preprocess) {前処理: Mel Spectrogram + Data Augmentation};

% === モデル層 ===
\node[modelblock, below=1.5cm of preprocess, xshift=-2.5cm] (student) {モデル: Student};
\node[modelblock, below=1.5cm of preprocess, xshift=2.5cm] (teacher) {モデル: Teacher};

% モデルをまとめるボックス
\begin{scope}[on background layer]
    \node[loopbox, fill=orange!5, fit=(student)(teacher), inner sep=15pt] (modelbox) {};
\end{scope}
\node[looplabel, fill=orange!20] at (modelbox.north) {Mean Teacher};

% === 出力層 ===
\node[block=purple!15, below=1.2cm of student] (pred_s) {強・弱予測};
\node[block=purple!15, below=1.2cm of teacher] (pred_t) {強・弱予測};

% === 損失計算 ===
\node[lossblock, below=1.5cm of $(pred_s)!0.5!(pred_t)$] (loss) {損失計算};

% 損失の詳細
\node[block=red!10, below=0.8cm of loss, minimum width=10cm, minimum height=1.5cm] (lossdetail) {
    \begin{tabular}{ccc}
        BCE Loss & + & Consistency Loss \\
        (強/弱) & & (Student $\leftrightarrow$ Teacher, 強/弱)
    \end{tabular}
};

% === 更新 ===
\node[block=yellow!20, below=1.2cm of lossdetail, xshift=-2.5cm] (bp) {Backpropagation\\(Student更新)};
\node[block=yellow!20, below=1.2cm of lossdetail, xshift=2.5cm] (ema) {EMA\\(Teacher更新)};

% === 矢印 ===
% データ→前処理
\draw[arrow] (databox.south) -- (preprocess.north);

% 前処理→モデル
\draw[arrow] (preprocess.south) -- ++(0,-0.5) -| (student.north);
\draw[arrow] (preprocess.south) -- ++(0,-0.5) -| (teacher.north);

% モデル→予測
\draw[arrow] (student.south) -- (pred_s.north);
\draw[arrow] (teacher.south) -- (pred_t.north);

% 予測→損失
\draw[arrow] (pred_s.south) -- ++(0,-0.5) -| (loss.north west);
\draw[arrow] (pred_t.south) -- ++(0,-0.5) -| (loss.north east);

% 損失→詳細
\draw[arrow] (loss.south) -- (lossdetail.north);

% 詳細→更新
\draw[arrow] (lossdetail.south) -- ++(0,-0.4) -| (bp.north);
\draw[arrow] (lossdetail.south) -- ++(0,-0.4) -| (ema.north);

% 更新→モデル (フィードバックループ)
\draw[dashedarrow, blue!60] (bp.west) -- ++(-1,0) |- node[note, pos=0.25, left] {Adam} (student.west);
\draw[dashedarrow, orange!60] (ema.east) -- ++(1,0) |- node[note, pos=0.25, right] {指数移動平均} (teacher.east);

% === エポックループ表示 ===
\begin{scope}[on background layer]
    \node[loopbox, draw=black!40, fill=blue!3, 
          fit=(preprocess)(modelbox)(pred_s)(pred_t)(loss)(lossdetail)(bp)(ema), 
          inner xsep=25pt, inner ysep=15pt] (epochloop) {};
\end{scope}
\node[looplabel, fill=blue!15, anchor=north east] at (epochloop.north east) {×150 epochs};

\end{tikzpicture}

\end{document}
